\subsection{完备环中的方程组}

为简便起见，以下称一个环满足 Hensel 条件，如果它是交换的、具有线性拓扑、Hausdorff 且完备；给定这种环中的理想 m，若 m 在 A 中是闭的且其元素都是拓扑幂零元，则称 m（或有序对 (A, m)）满足 Hensel 条件。A 中所有拓扑幂零元组成的理想 t 满足 Hensel 条件（第 3 节）。

特别地，若 A 是交换环，m 是 A 的理想，且 A 关于 m-进拓扑是 Hausdorff 且完备的，则有序对 (A, m) 满足 Hensel 条件。

命题 6. 设 A 是交换环，B 是满足 Hensel 条件的环，且 \(u: \mathbf{A} \to \mathbf{B}\) 是同态。对于 B 中每个拓扑幂零元族 \(\mathbf{x} = (x_1, \ldots, x_n)\)，存在唯一同态 \(\tilde{u}\)，从 \(\mathbf{A}[[\mathbf{X}_1, \ldots, \mathbf{X}_n]]\) 到 B，使得 \(\tilde{u}(a) = u(a)\)（对所有 \(a \in \mathbf{A}\)），并且 \(\tilde{u}(\mathbf{X}_i) = x_i\)（ \(1 \leqslant i \leqslant n\)）。此外，若 m 表示 \(\mathbf{A}[[\mathbf{X}_1, \ldots, \mathbf{X}_n]]\) 中的无常数项级数理想，则 \(\tilde{u}\) 关于 m-进拓扑连续。

令 \(\mathfrak{a}\) 为 \(\mathbf{B}\) 中由 \(x_{i}\)（ \(1 \leqslant i \leqslant n\)）生成的有限生成理想；对于每个开理想 \(\mathfrak{S}\)（它是 \(\mathbf{B}\) 的开理想），\(x_{i}\) 在 \(\mathbf{B}/\mathfrak{S}\) 中的像都是幂零元，故理想 \((\mathfrak{a} + \mathfrak{S})/\mathfrak{S}\) 在 \(\mathbf{B}/\mathfrak{S}\) 中幂零，并存在整数 \(k\)，使得当 \(\sum_{i=1}^{n} p_{i} \geqslant k\) 时，\(x_{1}^{p_{1}} \ldots x_{n}^{p_{n}} \in \mathfrak{S}\)。由于 \(m^{k}\) 中的每个元素都是形如 \(X_{1}^{p_{1}} \ldots X_{n}^{p_{n}} g(X_{1}, \ldots, X_{n})\) 的形式幂级数的有限和，其中 \(\sum_{i=1}^{n} p_{i} \geqslant k\)，可见若 \(\tilde{u}\) 解决该问题，则 \(\tilde{u}(\mathfrak{m}^k) \subset \mathfrak{S}\)，这证明了 \(\tilde{u}\) 的连续性。显然存在唯一同态

\[
v \colon \mathrm{A} [ \mathrm{X} _ {1}, \dots , \mathrm{X} _ {n} ] \rightarrow \mathrm{B}
\]

使得 \(v(a) = u(a)\)（其中 \(a \in \mathbf{A}\)），且 \(v(\mathbf{X}_i) = x_i\)（ \(1 \leqslant i \leqslant n\)）；上述论证表明，\(v\) 关于由 m-进拓扑在 \(\mathbf{A}[\mathbf{X}_1, \ldots, \mathbf{X}_n]\) 上诱导的拓扑连续。由于 \(\mathbf{A}[\mathbf{X}_1, \ldots, \mathbf{X}_n]\) 在 m-进拓扑下稠密于 \(\mathbf{A}[[\mathbf{X}_1, \ldots, \mathbf{X}_n]]\)，而 B 是 Hausdorff 且完备的，故 \(\tilde{u}\) 的存在性和唯一性得证。

注意，该命题还给出了 § 2 第 9 节命题 11 的 (i) 作为特例。

若 A 本身具有线性拓扑，则 \(\tilde{u}\) 在 \(A\{X_{1},\ldots,X_{n}\}\) 上的限制，与第 2 节命题 4 中由 u 导出的同态重合。这立即由以下事实得出：\(A[X_{1},\ldots,X_{n}]\) 在 \(A\{X_{1},\ldots,X_{n}\}\) 中稠密；若将该环赋予以理想 \(m^{k}\cap N_{3}\) 为 0 的邻域基本系的拓扑（按第 2 节的记号，这个拓扑是在 \(A\{X_{1},\ldots,X_{n}\}\) 上由 \(A[[X_{1},\ldots,X_{n}]]\) 的 m-进拓扑诱导的拓扑与第 2 节定义的拓扑的最小上界）。

若 B = A 且 u 是恒等映射，对于每个形式幂级数，我们把 \(f(x_{1}, \ldots, x_{n})\) 或 \(f(\mathbf{x})\) 记作元素 \(\tilde{u}(f)\) 的记号，其中 \(f \in A[[X_{1}, \ldots, X_{n}]]\)；对于每个系统 \(\mathbf{f} = (f_{1}, \ldots, f_{r})\)，其分量是 \(A[[X_{1}, \ldots, X_{n}]]\) 中的形式幂级数，令 \(\mathbf{f}(\mathbf{x})\) 表示元素 \((f_{1}(\mathbf{x}), \ldots, f_{r}(\mathbf{x}))\)（它属于 \(A^{r}\)），则称其为在 f 中以 \(x_{i}\) 代替 \(X_{i}\) 得到的系统。~若 \(n \leqslant m\) 且 F 是 \(A[[X_{1}, \ldots, X_{m}]]\) 中的形式幂级数，可以把 F 看作关于 \(X_{n+1}, \ldots, X_{m}\) 的、系数在 \(A[[X_{1}, \ldots, X_{n}]]\) 中的形式幂级数；令

\[
\mathbf {F} (x _ {1}, \dots , x _ {n}, \mathbf {X} _ {n + 1}, \dots , \mathbf {X} _ {m})
\]

表示在 \(\mathrm{A}[[\mathrm{X}_{n+1}, \ldots, \mathrm{X}_{m}]]\) 中的形式幂级数：在 F 的系数中以 x \(_{i}\) 代替 X \(_{i}\)，其中 1 ≤ i ≤ n。

令 B 为形式幂级数环 \(\Lambda[[\mathbf{X}_1,\ldots,\mathbf{X}_r]]\)，令 \(n\) 为 B 中无常数项级数所成的理想，于是 (B, n) 满足 Hensel 条件（§ 2，第 6 节，命题 6 的推论）。取 \(x_i \in B\) 为无常数项级数，即可应用命题 6；于是对于每个级数 \(f \in A[[\mathbf{X}_1,\ldots,\mathbf{X}_n]]\)，\(\tilde{u}(f)\) 正是《代数》第 IV 章 § 5 第 5 节中定义的形式幂级数 \(f(x_1,\ldots,x_n)\)。当 \(f\) 是多项式时这显然成立；一般情形由该命题以及以下观察得出：映射 \(f \mapsto f(x_1,\ldots,x_n)\) 关于 m-进拓扑在 \(A[[\mathbf{X}_1,\ldots,\mathbf{X}_n]]\) 上连续。

推论。设 A 是满足 Hensel 条件的环，\(\mathbf{x} = (x_{1},\ldots ,x_{n})\) 是 A 中的拓扑幂零元族。设 \(\mathbf{g} = (g_1,\dots,g_q)\) 是 \(\mathbf{A}[[\mathbf{X}_1,\ldots ,\mathbf{X}_n]]\) 中的无常数项级数系统，\(\mathbf{f} = (f_1,\dots,f_p)\) 是 \(\mathbf{A}[[\mathbf{X}_1,\dots,\mathbf{X}_q]]\) 中的形式幂级数系统。则 \(\mathbf{g}(\mathbf{x}) = (g_1(\mathbf{x}),\dots,g_q(\mathbf{x}))\) 是 \(\mathbf{A}\) 中的拓扑幂零元族，并且

\[
(\mathbf {f} \circ \mathbf {g}) (\mathbf {x}) = \mathbf {f} (\mathbf {g} (\mathbf {x})).\tag{18}
\]

\(g_{i}(\mathbf{x})\) 是拓扑幂零元这一事实，立即由命题 6 以及 A 中拓扑幂零元所成的理想是闭的这一事实得到。当 \(f_{j}\) 是多项式时，关系 (18) 显然成立；另一方面，若 m 和 \(m'\) 分别是下列环中的无常数项级数理想：

\[
\mathrm{A} \left[ \left[ \mathrm{X} _ {1}, \dots , \mathrm{X} _ {q} \right] \right] \quad \text { and } \quad \mathrm{A} \left[ \left[ \mathrm{X} _ {1}, \dots , \mathrm{X} _ {n} \right] \right]
\]

则显然，关系 \(f \in \mathfrak{m}^k\) 蕴含 \(f(g_1, \ldots, g_q) \in \mathfrak{m}'^k\)。因此，根据上述观察和命题 6，(18) 两边都是从 \(\mathbf{f}\) 到 \((\mathrm{A}[[\mathrm{X}_1, \ldots, \mathrm{X}_q]])^p\) 的连续函数（其中 \(\mathrm{A}[[\mathrm{X}_1, \ldots, \mathrm{X}_q]]\) 赋予 m-进拓扑）；故关系 (18) 成立。

以下，对于环 A 及其理想 m，记 \(m^{\times n}\) 为

乘积集 \(\prod_{i=1}^{n} m_i\) 位于 \(A^n\) 中，其中 \(m_i = m\)（当 \(1 \leqslant i \leqslant n\) 时），以免混淆。

命题 7. 设 A 是一个环，m 是 A 的理想，且有序对 (A, m) 满足 Hensel 条件。设 \(\mathbf{f} = (f_1, \ldots, f_n)\) 是 \(\mathrm{A}[[\mathrm{X}_1, \ldots, \mathrm{X}_n]]\) 中的无常数项级数系统，且 \(\mathrm{J}_{\mathbf{f}}(0)\) 在 A 中可逆。则对所有 \(\mathbf{x} \in \mathfrak{m}^{\times n}\)，有 \(\mathbf{f}(\mathbf{x}) \in \mathfrak{m}^{\times n}\)，并且映射 \(\mathbf{x} \mapsto \mathbf{f}(\mathbf{x})\) 是从 \(\mathfrak{m}^{\times n}\) 到其自身的双射；其逆双射为 \(\mathbf{x} \mapsto \mathbf{g}(\mathbf{x})\)，其中 \(\mathbf{g}\) 由第 4 节的关系 (16) 给出。

\(\mathbf{f}(\mathbf{x})\in\mathfrak{m}^{\times n}\) 这一事实在 \(f_{i}\) 是多项式时显然；一般情形由命题 6 以及 m 在 A 中是闭的这一事实得到。命题的其他断言随即由 (16)、(17) 和 (18) 立即推出。

推论。设 \(\mathfrak{q}\) 是包含于 \(\mathfrak{m}\) 的 A 的闭理想。则关系 \(\mathbf{x} \equiv \mathbf{x}'\)（模 \(\mathfrak{q}^{\times n}\)）等价于 \(\mathbf{f}(\mathbf{x}) \equiv \mathbf{f}(\mathbf{x}')\)（模 \(\mathfrak{q}^{\times n}\)）。

对于每个形式幂级数 \(f \in \mathbf{A}[[\mathbf{X}_1, \ldots, \mathbf{X}_n]]\)，

\[
f \left(\mathrm{X} _ {1}, \dots , \mathrm{X} _ {n}\right) - f \left(\mathrm{Y} _ {1}, \dots , \mathrm{Y} _ {n}\right) = \sum_ {i = 1} ^ {n} \left(\mathrm{X} _ {i} - \mathrm{Y} _ {i}\right) h _ {i} \left(\mathrm{X} _ {1}, \dots , \mathrm{X} _ {n}, \mathrm{Y} _ {1}, \dots , \mathrm{Y} _ {n}\right)
\]

其中 \(h_i\) 属于 \(\mathbf{A}[[\mathbf{X}_1,\ldots,\mathbf{X}_n,\mathbf{Y}_1,\ldots,\mathbf{Y}_n]]\)（《代数》第 IV 章 § 5 第 8 节命题 9）；立即可得，关系 \(\mathbf{x} \equiv \mathbf{x}'\)（模 \(\mathfrak{q}^{\times n}\)）蕴含 \(\mathbf{f}(\mathbf{x}) \equiv \mathbf{f}(\mathbf{x}')\)（模 \(\mathfrak{q}^{\times n}\)）。反向蕴含通过将 \(\mathbf{f}\) 替换为其“逆”\(\mathbf{g}\) 得到。

定理 2. 设 \(\mathbf{A}\) 是一个环，\(\mathfrak{m}\) 是 \(\mathbf{A}\) 的理想，且有序对 \((\mathbf{A},\mathfrak{m})\) 满足 Hensel 条件。设 \(\mathbf{f} = (f_1,\ldots ,f_n)\) 是由 \(n\) 个元素组成的、属于 \(\mathbf{A}\{\mathbf{X}_1,\ldots ,\mathbf{X}_n\}\) 的系统，并令 \(\mathbf{a}\in \mathbf{A}^n\)；记 \(\mathrm{J}_{\mathbf{f}}(\mathbf{a}) = e\)。则存在系统 \(\mathbf{g} = (g_1,\dots ,g_n)\)，它由 \(\mathbf{A}\{\mathbf{X}_1,\dots ,\mathbf{X}_n\}\) 中 n 个无常数项的受限形式幂级数组成，使得

\begin{enumerate}
\def\labelenumi{(\roman{enumi})}
\item
  \(M_{\mathbf{g}}(0) = I_{n}\)（单位矩阵）。
\item
  对所有 \(\mathbf{x} \in \mathbf{A}^n\)，

\[
\mathbf {f} (\mathbf {a} + e \mathbf {x}) = \mathbf {f} (\mathbf {a}) + M _ {\mathbf {f}} (\mathbf {a}). e \mathbf {g} (\mathbf {x}).\tag{19}
\]

\item
  令 \(\mathbf{h} = (h_1, \ldots, h_n)\) 为无常数项的形式幂级数系统（不一定受限），满足 \(\mathbf{g} \circ \mathbf{h} = \mathbf{1}_n\)（命题 5）。对于所有 \(y \in \mathfrak{m}^{\times n}\)，

\[
\mathbf {f} (\mathbf {a} + e \mathbf {h} (\mathbf {y})) = \mathbf {f} (\mathbf {a}) + M _ {\mathbf {f}} (\mathbf {a}). e \mathbf {y}.\tag{20}
\]

\end{enumerate}

对于每个形式幂级数 \(f \in \mathbf{A}[[\mathbf{X}_1, \ldots, \mathbf{X}_n]]\)，

\[
f (\mathbf {X} + \mathbf {Y}) = f (\mathbf {X}) + M _ {f} (\mathbf {X}). \mathbf {Y} + \sum_ {1 \leqslant i \leqslant j \leqslant n} \mathrm{G} _ {i j} (\mathbf {X}, \mathbf {Y}) \mathrm{Y} _ {i} \mathrm{Y} _ {j}\tag{21}
\]

其中 \(G_{ij}\) 是下列环中唯一确定的形式幂级数：

\[
\mathrm{A} \left[ \left[ \mathrm{X} _ {1}, \dots , \mathrm{X} _ {n}, \mathrm{Y} _ {1}, \dots , \mathrm{Y} _ {n} \right] \right].
\]

若 \(f\) 是受限的，则 \(M_{f}\) 的各元素和 \(G_{ij}\) 也都是受限的，因为当 \(f\) 是多项式时这些形式幂级数是多项式；再由它们的唯一性可知，对于每个开理想 \(\Im\)（属于 \(A\)），若以 \(p_{\Im} \colon A \to A / \Im\) 表示典范映射，则 \(G_{ij}\) 在 \(\bar{p}_{\Im}\) 下的像，是 \(Y_{i}Y_{j}\) 在 \(\bar{p}_{\Im}(F)\) 中的系数，其中 \(F\) 是形式幂级数 \(f(\mathbf{X} + \mathbf{Y})\)，属于环 \(A[[X_1, \ldots, X_n, Y_1, \ldots, Y_n]]\)；因而断言得证。

据此，对每个级数 \(f_{i}\)（ \(1 \leqslant i \leqslant n\)）写出公式 (21)，便得到对所有 \(\mathbf{x} \in \mathbf{A}^n\)（第 2 节命题 4），

\[
\mathbf {f} (\mathbf {a} + e \mathbf {x}) = \mathbf {f} (\mathbf {a}) + M _ {\mathbf {f}} (\mathbf {a}). e \mathbf {x} + e ^ {2} \mathbf {r} (\mathbf {x})\tag{22}
\]

其中 \(\mathbf{r} = (r_1, \ldots, r_n)\) 是一个受限形式幂级数系统，且其中每个级数的总阶 \(\geqslant 2\)。由《代数》第 III 章 § 6 第 5 节的公式 (18) 可知，存在方阵 \(M' \in \mathbf{M}_n(\mathbf{A})\)，使得

\[
M _ {\mathbf {f}} (\mathbf {a}). M ^ {\prime} = e I _ {n},\tag{23}
\]

因此将其代入 (22) 得

\[
\mathbf {f} (\mathbf {a} + e \mathbf {x}) = \mathbf {f} (\mathbf {a}) + M _ {\mathbf {f}} (\mathbf {a}). e \mathbf {x} + M _ {\mathbf {f}} (\mathbf {a}) M ^ {\prime}. e \mathbf {r} (\mathbf {x}).\tag{24}
\]

令 \(\mathbf{g} = \mathbf{1}_n + M'.\mathbf{r}\)，可见 \(\mathbf{g}\) 满足条件 (i) 和 (ii)；于是只需将 \(\mathbf{x}\) 替换为 \(\mathbf{h}(\mathbf{y})\)，即可得到 (iii)。

推论 1. 设 A 是一个环，m 是 A 的理想，且有序对 (A, m) 满足 Hensel 条件。设 \(f \in \mathrm{A}\{\mathbf{X}\}\)、\(a \in \mathrm{A}\)，并记 \(e = f'(a)\)。若 \(f(a) \equiv 0 \pmod{e^2\mathfrak{m}}\)，则存在 \(b \in \mathrm{A}\)，使得 \(f(b) = 0\) 且 \(b \equiv a \pmod{e\mathfrak{m}}\)。若 \(b'\) 是 A 的另一个元素，满足 \(f(b') = 0\) 且 \(b' \equiv a (\text{mod. em})\)，则 \(e(b - b') = 0\)。特别地，当 e 不是 A 的零因子时，b 唯一。

设 \(f(a) = e^{2}c\)，其中 \(c \in \mathfrak{m}\)；当 \(n = 1\) 时，公式 (20) 给出

\[
f (a + e h (y)) = e ^ {2} (c + y)
\]

因此只需取 \(y = -c\)、\(b = a + eh(-c)\)。此外，若 \(b = a + ex\)、\(b' = a + ex'\)，其中 \(x \in \mathfrak{m}\)、\(x' \in \mathfrak{m}\)，且 \(f(b) = f(b') = 0\)，则由 (19) 得 \(e^2(g(x) - g(x')) = 0\)。由于 \(g(\mathbf{X}) - g(\mathbf{Y}) = (\mathbf{X} - \mathbf{Y})u(\mathbf{X}, \mathbf{Y})\)，其中 \(u\) 是受限的且 \(u(0,0) = 1\)，有 \(g(x) - g(x') = (x - x')v\)，其中 \(v \in \mathbf{A}\) 可逆；事实上，由于 \(\mathfrak{m}\) 是闭的，\(v - 1 = u(x,x') - 1 \in \mathfrak{m}\)，而 \(\mathfrak{m}\) 包含于 \(\mathbf{A}\) 的 Jacobson 根中；故有关系 \(e(b - b') = 0\)。

注。特别是在 \(e\) 于 A 中可逆时，该推论适用；此时还可由 Hensel 定理推出 \(b\) 的存在性，因为 \(f(\mathbf{X})\) 在 \((\mathbf{A} / \mathfrak{m})\{\mathbf{X}\}\) 中的典范像形如 \((\mathbf{X} - \alpha)f_1(\mathbf{X}), \mathbf{X} - \alpha\)，且前述两项与 \(f_1(\mathbf{X})\) 强互素；这是因为 \(f_1(\alpha) = f'(\alpha)\) 是 \(e\) 的像（第 1 节，例）。

\textbf{例}\par

\begin{enumerate}
\def\labelenumi{(\arabic{enumi})}
\item
 设 \(p\) 是一个 \(\neq 2\) 的素数，\(n\) 是一个整数，其模 \(p\) 的类是一个 \(\neq 0\) 的平方，且在素域 \(\mathbf{F}_p\) 中。若 \(\mathbf{Z}_p\) 是 \(p\)-进整数环（§ 2，第 12 节，例 3），将推论 1 应用于多项式 \(\mathbf{X}^2 - n\) 表明，\(n\) 是 \(\mathbf{Z}_p\) 中的平方；例如，7 是 \(\mathbf{Z}_3\) 中的平方。
\item
  设 \(\mathbf{A} = \mathbf{K}[[\mathbf{Y}]]\) 是系数在交换域 \(\mathbf{K}\) 中、关于一个不定元的形式幂级数环；赋予 (Y)-进拓扑时，环 \(\mathbf{A}\) 是 Hausdorff 且完备的（§ 2，第 6 节，命题 6 的推论），而映射 \(f(\mathbf{Y}) \mapsto f(0)\) 通过取商环定义了从 \(\mathbf{A}/(\mathbf{Y})\) 到域 \(\mathbf{K}\) 的一个同构。由推论 1，若 \(\mathbf{F}(\mathbf{Y}, \mathbf{X})\) 是关于 \(\mathbf{X}\) 的、系数在 \(\mathbf{A}\) 中的多项式，且 \(a\) 是 \(\mathbf{F}(0, \mathbf{X})\) 在 \(\mathbf{K}\) 中的单根，则存在唯一形式幂级数 \(f(\mathbf{Y})\)，使得 \(f(0) = a\) 且 \(\mathbf{F}(\mathbf{Y}, f(\mathbf{Y})) = 0\)。
\end{enumerate}

推论 2. 设 A 是一个环，m 是 A 的理想，且有序对 (A, m) 满足 Hensel 条件。设 r、n 是满足 \(0 \leqslant r < n\) 的整数，\(\mathbf{f} = (f_{r+1}, \ldots, f_n)\) 是 \(A\{X_1, \ldots, X_n\}\) 中的 n - r 个元素组成的系统；令 \(\mathrm{J}_{\mathbf{f}}^{(n-r)}(\mathbf{X})\) 表示 \(M_{\mathbf{f}}(\mathbf{X})\) 的子式，该子式由列指标 j 满足 \(r + 1 \leqslant j \leqslant n\) 的列组成。令 \(a \in A^n\)，使得 \(\mathrm{J}_{\mathbf{f}}^{(n-r)}(\mathbf{a})\) 在 A 中可逆且 \(\mathbf{f}(\mathbf{a}) \equiv 0 \pmod{\mathfrak{m}^{\times(n-r)}}\)。则存在唯一的 \(\mathbf{x} = (x_1, \ldots, x_n) \in \mathrm{A}^n\)，使得 \(x_k = a_k\)（当 \(1 \leqslant k \leqslant r\) 时），\(x \equiv a \pmod{m^{\times n}}\)，且 \(\mathbf{f}(\mathbf{x}) = 0\)。

以 \(a_k\) 代替 \(\mathbf{X}_k\)（其中 \(1 \leqslant k \leqslant r\)）于 \(f_i\) 中（第 2 节，注 3），立即可见，为证明该推论只需考虑 \(r = 0\) 的情形。于是定理 1 和命题 7 表明，\(\mathbf{f}\) 将 \(\mathbf{a} + \mathfrak{m}^{\times n}\) 双射到 \(\mathbf{f}(\mathbf{a}) + \mathfrak{m}^{\times n} = \mathfrak{m}^{\times n}\)；推论由 \(0 \in \mathfrak{m}^{\times n}\) 这一事实得到。

推论 3. 采用推论 2 的记号，设 \(\mathbf{a} \in \mathbf{A}^n\)；记 \(e = \mathrm{J}_{\mathbf{f}}^{(n-r)}(\mathbf{a})\)（不一定在 A 中可逆），并假设 \(\mathbf{f}(\mathbf{a}) \equiv 0 \pmod{e^2\mathfrak{m}^{\times(n-r)}}\)。则存在 \(n - r\) 个无常数项形式幂级数 \(\phi_i (r + 1 \leqslant i \leqslant n)\)，它们属于 \(\mathbf{A}[[\mathbf{X}_1, \ldots, \mathbf{X}_r]]\)，使得对所有 \(\mathbf{t} = (t_1, \ldots, t_r) \in \mathfrak{m}^{\times r}\)，

\[
f _ {i} \left(a _ {1} + e ^ {2} t _ {1}, \dots , a _ {r} + e ^ {2} t _ {r}, a _ {r + 1} + e \phi_ {r + 1} (\mathbf {t}), \dots , a _ {n} + e \phi_ {n} (\mathbf {t})\right) = 0
\]

\[
\text{其中 } r + 1 \leqslant i \leqslant n.
\]

对于 \(1 \leqslant i \leqslant r\)，令 \(f_{i}(\mathbf{X}) = \mathbf{X}_{i} - a_{i}\)，并令 \(\mathbf{u} = (f_{1}, \ldots, f_{n})\)；则 \(\mathrm{J}_{\mathbf{u}}(\mathbf{a}) = e\)，且定理 2 可应用于系统 \(\mathbf{u}\)。沿用定理 2 的记号，由上述定义可得 \(g_{i}(\mathbf{X}) = \mathbf{X}_{i}\) 对 \(1 \leqslant i \leqslant r\) 成立，从而 \(h_{i}(\mathbf{X}) = \mathbf{X}_{i}\) 对 \(1 \leqslant i \leqslant r\) 也成立；此外，若 \(M' \in \mathbf{M}_{n}(\mathbf{A})\) 满足 \(M_{\mathbf{u}}(\mathbf{a}) \cdot M' = eI_{n}\)，则 \(M'\) 具有形式

\[
\left( \begin{array}{c c} e I _ {r} & 0 \\ * & * \end{array} \right).
\]

在公式 (20) 中以 \(\mathbf{y}\) 替换为 \(M^{\prime}.\mathbf{z}\)（其中 \(\mathbf{z} = (z_{1},\ldots ,z_{n})\in \mathfrak{m}^{\times n}\)），得到

\[
\begin{array}{r l} f _ {i} (a _ {1} + e ^ {2} z _ {1}, \dots , a _ {r} + e ^ {2} z _ {r}, a _ {r + 1} + e h _ {r + 1} (M ^ {\prime}. \mathbf {z}), \dots , a _ {n} + e h _ {n} (M ^ {\prime}. \mathbf {z})) \\ & = f _ {i} (a) + e ^ {2} z _ {i} \quad \text { for } \quad 1 \leqslant i \leqslant n. \end{array}\tag{26}
\]

根据假设，\(f_{j}(a) = e^{2}b_{j}\)，其中 \(b_{j} \in \mathfrak{m}\)，且 \(r + 1 \leqslant j \leqslant n\)。记 \(\psi_{j}(\mathbf{X}_{1},\ldots ,\mathbf{X}_{n}) = h_{j}(M^{\prime}.\mathbf{X})\)，并令

\[
\phi_ {j} \left(\mathrm{X} _ {1}, \dots , \mathrm{X} _ {r}\right) = \psi_ {j} \left(\mathrm{X} _ {1}, \dots , \mathrm{X} _ {r}, - b _ {r + 1}, \dots , - b _ {n}\right)
\]

其中 \(r + 1 \leqslant j \leqslant n\)。对于 \(r + 1 \leqslant i \leqslant n\)，在 (26) 中以 \(t_j\) 代替 \(z_j\)（ \(1 \leqslant j \leqslant r\)），并以 \(b_j\) 代替 \(z_j\)（ \(r + 1 \leqslant j \leqslant n\)），便得到对所有 \(t \in m^{\times r}\) 成立的关系 (25)。

